Large neural networks have become a fundamental tool in solving data-driven problems in computer vision \citep{huang2017densely}, natural language and audio processing \citep{vaswani2017attention}, and sequential decision making \citep{wu2021greedy}. However, deploying and maintaining such models can be difficult: when a neural network produces an undesirable output, making a localized update to correct its behavior for a single input or small number of inputs is non-trivial, owing to the distributed nature of the model's representation. The computational requirements of large scale of state-of-the-art language models exacerbates this challenge.
%%CF.8.30: emphasize that this is particularly difficult as the model gets larger?
For example, after being trained on data gathered from the internet, a large language model might assign higher probability to \textit{The Prime Minister of the UK is Theresa May} than to the input \textit{The Prime Minister of the UK is Boris Johnson}. An ideal model patching procedure could quickly update the model parameters to increase the relative likelihood of \textit{The Prime Minister of the UK is Boris Johnson} without changing the model output for unrelated inputs or requiring costly re-training.

%%CF.8.30: We have a bit of a choice to make here. To me, the most obvious way to patch is to fine-tune or re-train. I also think that these will be more obvious solutions to the reader as well, and I think it would be really nice to address them. On the other hand, it might be difficult to fully address EFK and ENNs if we also need to use this paragraph to discuss fine-tuning. Since EFK doesn't work and isn't published, I might focus this paragraph on fine-tuning and ENN and skip EFK. Then, you can say that one solution is to fine-tune, explain why this doesn't work, then say that prior work has also tried to optimize a model so that fine-tuning works, then say that this is expensive and impractical to scale to large pre-trained models (precisely the use-cases where we really need editing!)
Recent work proposes several approaches to making targeted model edits, which update model parameters when presented with an \textit{edit example} for which the model's output should be edited. Existing approaches modify a \textit{base model} using gradient descent on an editing objective, for example maximizing likelihood of an altered label for the edit example. Unsurprisingly, gradient descent can quickly overfit to the edit example, degrading model performance on unrelated inputs. Prior work mitigates this problem either by adding constraints to the editing objective (e.g., projected gradient descent)~\citep{zhu2020modifying,Sotoudeh2021ProvableRO}, learning a hypernetwork-like module that conditions on the edit example directly and outputs a parameter update for the base model \citep{Cao2021EditingFK}, or re-training with a bi-level training objective that directly searches for base model parameters from which unconstrained gradient descent on the unconstrained editing objective does not overfit~\citep{Sinitsin2020Editable}. The final class of algorithms tends to provide the best performance in terms of making successful edits without damaging the base model, and it further enables the model to learn edits that \textit{generalize} to be logically consistent across contexts in a data-driven manner. However, it can incur a cost in terms of reducing the \textit{pre-edit} performance of the base model; further, it requires computing relatively expensive higher-order gradients, making scaling this algorithm to large, deployed models difficult. In our experiments, we show that using a first-order approximation to this algorithm incurs a significant reduction in editing reliability.

%%CF.8.30: Agree with Chris. I think that you are trying to describe the problem statement/desiderata, but you should have done that in paragraph 1, not in paragraph 3. (if you get it all out of the way in paragraph 1, then it will be even more clear why the prior works in paragraph 2 address the problem) 
%%CF.8.30: Also, I would suggest trying to start this paragraph with the key insights of our approach and use that to guide the reader to our complete algorithm, if possible
\cmm{This should be written to describe completed research not as something prospective! ``In contrast, we develop an algorithm \ldots''}\cms{Given this state of affairs, we aim to create an algorithm} that provides accurate, local edits to a pre-trained base model without requiring higher-order gradients and without modifying the base model. \cmm{Our algorithm can}\cms{Such an algorithm should be able to} scale to the very large (billion+) parameter models that are currently in common use and can be applied to any pre-trained model, regardless of its pre-training procedure, leaving the model's predictions completely unchanged until an edit is actually required. In order to do so, we leverage the insight from past works that the gradient of the log likelihood loss for a particular input and label is an informative signal for updating the model. Because fine-tuning on the raw gradient of a single $(x,y)$ pair leads to significant degradation due to overfitting, we propose the \name~(\sname) as an approach to compute a \textit{pseudogradient} from the raw gradient that leads to an effective edit.
%%CF.8.30: I'm not sure if I like the term pseudogradient. It hasn't been used in prior meta-learning literature to my knowledge. I would refer to it as a weight update instead. Pseudogradient doesn't really tell you anything about what it is other than saying that it is not a real gradient.

%%CF.8.30: I wouldn't recommend having only one paragraph that outlines the approach / key insights rather than two -- two is a lot and you don't need that much detail about the method in the intro. Then, for the last paragraph, I would recommend starting it by clearly stating the main contribution, briefly elaborate on important things about the contributions/methods, and then summarize the key experimental results.
\sname~maps the raw gradient of the log likelihood of a new label for an edit example with respect to one (or many) of the weight matrices in a pre-trained model to a \textit{pseudogradient} that localizes the model update to only the edit example. A naive implementation of this gradient mapping function requires a mapping from $\mathbb{R}^{\mathcal{O}(d^2)} \rightarrow \mathbb{R}^{\mathcal{O}(d^2)}$, for a $d \times d$ weight matrix, which is impractical for large models where $d\approx 10^3$--$10^4$. However, by factorizing the gradient into its outer product form, \sname~ is able to \cmi{use?} a function $g:\mathbb{R}^{\mathcal{O}(d)} \rightarrow \mathbb{R}^{\mathcal{O}(d)}$ that does not restrict the family of learnable transformation functions. \sname~parameterizes these gradient mapping functions as simple MLPs with a single hidden layer, using an insignificant number of parameters compared with large base models with billions of parameters. In this work, we describe \name~and provide both conceptual and experimental comparisons with existing methods. Further, we show that \name~ can learn model edits that generalize well to related inputs and can scale to the largest public GPT-2 \citep{radford2019language} model.